\originalchapter{8.334作业3：标度、耦合场与长程RG}
\psource{Mehran Kardar，8.334 Spring 2014，PS3，截止2014年3月21日。}
\begin{exercise}[原题1(a)--(c)]\label{p:334-3-1}
二维$n$分量非线性$\sigma$模型给$\dot T=(n-2)T^2/(2\pi)$、$\dot h=2h$，点为对$\ell$求导。(a)$T\to0$相关长；(b)$T,h\to0$奇异自由能；(c)$h=0$磁化率。
\end{exercise}
\begin{solution}
先假设$n>2$。积分$1/T(\ell)=1/T_0-(n-2)\ell/(2\pi)$，到$T(\ell_*)=O(1)$为止，$\ell_*=2\pi/[(n-2)T_0]+O(1)$，$\xi\sim a\e^{\ell_*}$。面积自由能$f_s=\e^{-2\ell_*}\Phi(h\e^{2\ell_*})$，即$f_s\sim\xi^{-2}\Phi(h\xi^2)$（指标按题给截断流，忽略高阶场异常维）。二次场导数$\chi\sim\xi^2\sim\exp[4\pi/((n-2)T)]$。

$n=2$的该微扰流恒0，不能除以$n-2$；XY需要涡旋的BKT方程。$n<2$亦不适用此渐近自由指数尺度，原公式没有无条件覆盖所有$n$。
\end{solution}
\begin{exercise}[原题2(a)--(e)]\label{p:334-3-2}
两标量$\beta H=\int[tm^2/2+K(\nabla m)^2/2-hm+L(\nabla^2\phi)^2/2+v\nabla m\cdot\nabla\phi]\dd^dx$。(a)Fourier式；(b)$q'=bq,m'(q')=m(q)/z,\phi'(q')=\phi(q)/y$的递推，忽略常数积分；(c)$K,L$固定点$y_t,y_h,y_v$；(d)$f_s=t^{2-\alpha}g(h/t^\Delta,v/t^\omega)$的指数；(e)$t,L$固定的另一点，哪些扰动相关。
\end{exercise}
\begin{solution}
$\beta H=\frac12\int_q(m(-q),\phi(-q))\left(\begin{smallmatrix}t+Kq^2&vq^2\\vq^2&Lq^4\end{smallmatrix}\right)\binom{m(q)}{\phi(q)}-hm(0)$。消去高动量Gaussian仅贡献自由能常数，低模改变量得$t'=b^{-d}z^2t$、$K'=b^{-d-2}z^2K$、$L'=b^{-d-4}y^2L$、$v'=b^{-d-2}zyv$、$h'=zh$。

取$z=b^{(d+2)/2}$、$y=b^{(d+4)/2}$，$K,L$固定，$y_t=2,y_h=(d+2)/2,y_v=1$。自由能体积缩放$b^{-d}$，选$b=t^{-1/2}$得$\alpha=2-d/2$、$\Delta=(d+2)/4$、$\omega=1/2$。$v$相关且把临界线移到$t=v^2/L$：配方有效$m$核$t-v^2/L+Kq^2$（$q\ne0$），稳定区$t>v^2/L$。不能在$t=0$保留有限$v$后继续使用不稳定Gaussian。

若固定$t,L$，则$z=b^{d/2}$、$y=b^{(d+4)/2}$，$K'=b^{-2}K$无关、$v'=v$边缘、$h'=b^{d/2}h$相关；这是有质量$m$配以弯曲无质量$\phi$的固定描述，$v$只在$t>v^2/L$的稳定区改变有效系数。两套场重标度反映选择不同主导核，不能混用其$y_t$。
\end{solution}
\begin{exercise}[原题3(a)--(j)]\label{p:334-3-3}
$d$维分为$n$个平行、$d-n$个垂直方向，$\beta H_0=\int[K(\nabla_\parallel m)^2/2+L(\nabla_\perp^2m)^2/2+tm^2/2-hm]\dd^nx_\parallel\dd^{d-n}x_\perp$，扰动$U=u\int m^4$。(a)Fourier；(b)$q'_\parallel=bq_\parallel,q'_\perp=cq_\perp,m'=m/z$递推；(c)$K,L$固定的$c,z,y_t,y_h$；(d)自由能、$\alpha,\Delta$；(e)零场$\chi$；(f)Gaussian两点和$\chi_0(q)$；(g)$U$的Fourier表达；(h)上临界维；(i)一阶展开用Wick化简；(j)$\chi(q)$与一阶$t_c$，无需积分。
\end{exercise}
\begin{solution}
记$D(q)=t+Kq_\parallel^2+Lq_\perp^4$，$H_0=\tfrac12\int_qD|m|^2-hm(0)$。设体积因子$B=b^nc^{d-n}$，有$K'=z^2K/(Bb^2)$、$L'=z^2L/(Bc^4)$、$t'=z^2t/B$、$h'=zh$。两梯度同时固定要求$c=b^{1/2}$，定义有效缩放维$D_e=n+(d-n)/2=(d+n)/2$，$z=b^{(D_e+2)/2}$。因此$y_t=2,y_h=(D_e+2)/2$，$f_s=b^{-D_e}f_s(b^2t,b^{y_h}h)$，$2-\alpha=D_e/2$、$\Delta=(D_e+2)/4$、$\chi=t^{-1}$。

$\avg{m(q)m(q')}_0=(2\pi)^d\delta(q+q')/D(q)$，$\chi_0=1/D$。$U=u\int_{q_1q_2q_3q_4}(2\pi)^d\delta(\sum q_i)\prod m(q_i)$。其递推$u'=uz^4/B^3=b^{4-D_e}u$，故$d_u=8-n$（在题给固定$n$的方向划分下）。

连接一阶展开$\avg{m_qm_{q'}}=\avg{m_qm_{q'}}_0-\avg{m_qm_{q'}U}_{0,c}$。顶点四条腿选两条接外腿有$4\times3=12$，另两条自收缩$I_1=\int_kD(k)^{-1}$，故$\chi(q)=D(q)^{-1}-12uI_1D(q)^{-2}+O(u^2)$。重写为$[t+12uI_1+Kq_\parallel^2+Lq_\perp^4]^{-1}$到一阶，$t_c=-12u\int_k[Kk_\parallel^2+Lk_\perp^4]^{-1}+O(u^2)$，需同一UV截断与IR可积条件；低维发散正是Gaussian失效信号。平行$\xi\sim t^{-1/2}$、垂直$\xi\sim t^{-1/4}$，同一个长度指数不能描述两方向。
\end{solution}
\begin{exercise}[原题4(a)--(g)]\label{p:334-3-4}
$N_s$分量自旋有$J(r)\propto r^{-d-\sigma}$长程项和短程LG项。(a)导出$t+K_2q^2+K_\sigma q^\sigma$及$u(m^2)^2$核；(b)$u=0,\sigma>2$递推与固定核；(c)$\sigma<2$的场重标度与$K_2$无关；(d)Gaussian$\nu,\eta,\gamma$、$d>2\sigma$时$u$；(e)微扰RG；(f)$d<2\sigma$的首阶指数，原题记$\epsilon=d-2\sigma$；(g)$J(r)\propto\e^{-r/a}$时行为。
\end{exercise}
\begin{solution}
长程非解析项的推导见习题\ref{p:334-2-3}；Fourier四次项$u\int_{q_1q_2q_3}(m_{q_1}\cdot m_{q_2})(m_{q_3}\cdot m_{-q_1-q_2-q_3})$。固定$q^2$核时$t'=b^2t$、$K_\sigma'=b^{2-\sigma}K_\sigma$，$\sigma>2$流向短程Gaussian。固定$q^\sigma$则Fourier场因子$z=b^{(d+\sigma)/2}$，$t'=b^\sigma t$、$K_2'=b^{\sigma-2}K_2$。于是$\nu=1/\sigma$，$\chi(q)\sim q^{-\sigma}=q^{-2+\eta}$给$\eta=2-\sigma$，$\gamma=1$，$u$维$2\sigma-d$。

令$K_\sigma=1$、$A_d=S_{d-1}/(2\pi)^d$，窄壳积分$I_1=A_d\Lambda^{d-\sigma}\dd\ell$、$I_2=A_d\Lambda^{d-2\sigma}\dd\ell$。为核对原题(e)的系数，将$m=\varphi+\chi$，$u(m^2)^2$中的两快腿部分写作$V_2=\int\chi_aB_{ab}\chi_b$，$B_{ab}=2u\varphi^2\delta_{ab}+4u\varphi_a\varphi_b$。于是$\operatorname{tr}B=2u(N_s+2)\varphi^2$，一次收缩给$\delta t=4(N_s+2)uI_1$；两顶角连接收缩给$-I_2\operatorname{tr}B^2=-4(N_s+8)u^2I_2(\varphi^2)^2$。展开$I_1(t)=A_d\Lambda^{d-\sigma}\dd\ell[1-t/\Lambda^\sigma+\cdots]$并缩放，对$u(m^2)^2$约定（没有$4!$）得到
\begin{align}
\dot t&=\sigma t+4(N_s+2)u A_d\Lambda^{d-\sigma}-4(N_s+2)ut A_d\Lambda^{d-2\sigma}+\cdots,\\
\dot u&=(2\sigma-d)u-4(N_s+8)u^2 A_d\Lambda^{d-2\sigma}+\cdots.
\end{align}
动量非解析项不由局部壳积分重新生成，其系数这一阶无修正。取$\bar\epsilon=2\sigma-d=-\epsilon>0$，平移质量至临界流形后，$u_*A_d\Lambda^{-\bar\epsilon}=\bar\epsilon/[4(N_s+8)]$，
\begin{equation}
\nu^{-1}=\sigma-\frac{N_s+2}{N_s+8}\bar\epsilon+O(\bar\epsilon^2),\qquad \eta=2-\sigma.
\end{equation}
$\gamma=\sigma\nu$在该长程固定点。注意原题$\epsilon$是负数，不可直接套通常正的$4-d$公式。有限程指数尾所有低阶矩有限，Fourier在$q=0$解析，流向短程普适类（非Gaussian时是短程Wilson--Fisher）。长程固定点与短程点的交界需比较$\sigma$与$2-\eta_{SR}$，首阶Gaussian分界只是近似。
\end{solution}
